% sections/analysis.tex -- analysis aligned with the subquestions (shared).
Analysis is aligned with the subquestions. For SQ1, I will analyse selected
workshop episodes multimodally, adapting the video microethnography
\citet{mcarthur2015} used for interface interaction, and read interfaces,
terms of service and model documentation as cultural documents. For SQ2,
reflexive thematic analysis of the focus groups \citep{braun2019} is read
alongside field notes, episodes and group-created rules. For SQ3, version
analysis charts each state of the work against the contributors who shaped
it, after the version timelines of \citet[10]{hodson2019}. A system log records events, inputs and
transformations; it does not independently establish intention, authorship,
recognition or responsibility. Contribution trajectories will therefore be
interpreted through participant accounts, negotiated attribution, contextual
observation, material versions and acknowledged uncertainty, not used as a
surveillance record. Discrepant cases, including refusal and non-use,
reflexive memos and a decision trail keep the interpretation accountable. Because models cannot be isolated from
facilitation, platform, group history or changing tool versions, claims will
be situated and analytically transferable rather than causal.
